\documentclass[11pt]{article}
\usepackage[a4paper,margin=24mm]{geometry}
\usepackage{amsmath,amssymb,amsthm,booktabs,hyperref}
\newtheorem{lemma}{Lemma}\newtheorem{theorem}{Theorem}\newtheorem{corollary}{Corollary}
\newcommand{\gen}[1]{\langle#1\rangle}\newcommand{\F}{\mathbb F}
\title{String separators, normal modules and transitive rank}
\author{A supplement to the complete Topic B common theory}
\date{12 October 2026}
\begin{document}\maketitle
Let $G$ be a finite group generated by involutions
$r_0,\ldots,r_{d-1}$, all distinct and nonidentity. It is a string C-group
if $[r_i,r_j]=1$ when $|i-j|>1$,
and with the intersection property
$\gen{r_i:i\in I}\cap\gen{r_j:j\in J}=\gen{r_i:i\in I\cap J}$.
In particular its generators, and every subfamily, are independent:
no generator belongs to the subgroup on the others. Its noncommuting
diagram is the subgraph of the labelled path having an edge precisely
when the adjacent generators fail to commute. An interval of labels
need not be a connected component of this graph.

This supplement uses the complete accepted Topic B common manuscript
\emph{Forest coordinates and relative generation},
not an unpaid application theorem. Its full sixteen pages remain an
explicit common input. We use its normal-quotient and normal-module
bounds, extremal-tree theorem, and the weak primitive independent-length
bound. The latter is needed for block images, which need not inherit IP.
We do not use the older general imprimitive irredundant-length theorem:
the normalized two-block image does inherit full IP, and strong induction
on the faithful permutation degree supplies its sharper bound.
Whiston's bound and equality clause, original Mar\'oti classification,
and the classical identifications stated there remain retained premises.
The separate Conder lower-order theorem is used only in the sharp
primitive section and is explicitly retained there.

The block normalization is borrowed from [CFLM14]. The new proof supplies the two-neighbor case omitted
from their two-large-block argument, a complete finite residual check,
and an explicit normal-module and cocycle treatment of equality.
It does not claim priority for their rank theorem or for the borrowed
normalization. Ordered string tuples and unordered marked sets are
kept distinct.

\section{Primitive quotients on a maximal block system}
Let a transitive string C-group on $km$ letters preserve a maximal
nontrivial block system with $m$ blocks of size $k$. Its block image
$P$ is primitive. Choose an inclusion-minimal distinguished subfamily
$L$ generating $P$, and put $C$ equal to the distinguished generators
outside $L$ commuting with every member of $L$. Denote the remaining
family by $R$.

If the noncommuting diagram of the images of $L$ is disconnected,
its commuting factors are normal in $P$. A primitive nontrivial normal
subgroup is transitive; two commuting transitive subgroups are regular.
Their intersection is central and trivial. Each regular normal factor
must be simple, since a proper characteristic or normal subgroup
would give a block system for the product. Thus $P=T\times T$ in the
left--right action on a nonabelian simple group $T$, and $m=|T|\geq60$.
There are exactly two diagram components. Each irredundant generating
family has length at most $\log_2|T|$, so $|L|\leq2\log_2m$.
There are at most four adjacent positions in $R$.

Otherwise the images of $L$ have a connected diagram. Therefore the
corresponding distinguished vertices occupy an interval and $|R|\leq2$.
Whiston gives $|L|\leq m-1$, with equality only for $P=S_m$.
The primitive block action is not cyclic of odd prime order, since
it is generated by involutions. Except when $m=2$, it consequently
has trivial centralizer. Thus $\gen C$ fixes each block. Its restriction
to one block is faithful: an element trivial on that block and commuting
with $\gen L$ is trivial on all blocks, since $L$ is transitive on them.
Hence $|C|\leq k-1$.

For the disconnected case this gives
\[
 d\leq2\log_2m+k+3<km/2\qquad(m\geq60,\ k\geq2).
\]
For example the difference $km/2-k-3-2\log_2m$ is increasing in $k$
and $m\geq60$, and is positive at $(k,m)=(2,60)$.

Suppose now $k,m>2$ and $L$ is connected. Initially $d\leq k+m$.
Equality would make $P=S_m$ and $\gen C$ a diagonal $S_k$ on the blocks.
The diagram of $C$ is connected because $S_k$ is centerless and directly
indecomposable. Both neighbors of the $L$ interval must be present;
all of $C$ lies on one side of it. The opposite neighbor centralizes
$C$. In the full block stabilizer, the centralizer of this diagonal
$S_k$ is precisely the group $S_m$ permuting the blocks consistently.
Since $L$ lies in this centralizer and maps onto $S_m$, it is the whole
centralizer. The opposite neighbor is therefore already in $\gen L$,
contrary to independence. Consequently
\[
 d\leq k+m-1\leq km/2\quad(k,m\geq3,\ km\geq10).
\]
The last inequality follows from $(k-2)(m-2)\geq2$; the only excluded
pair is $(3,3)$, which is part of the finite degree-nine boundary.

\section{Two-point blocks and the equality marking}
Assume $k=2$ and $m\geq5$. Write $V=\mathbb F_2^m$ for the base
module of $C_2\wr S_m$, $W=G\cap V$, and $P=G/W$.
Any member of $C$ is the constant flip $\lambda=(1,\ldots,1)$,
so $|C|\leq1$. For a connected $L$ we have $|R|\leq2$.
If $P\ne S_m,A_m$, the accepted primitive independent-length bound
$\mu'(P)\leq m-3$ gives $d\leq m$, apart from the primitive degree-six
copy of $S_5$.

The $A_m$-invariant subspaces of $V$, for $m\geq5$, are
$0,\gen\lambda,E,V$, where $E$ is the even-weight space (with the
usual coincident or nested spaces understood). Indeed, from any
nonconstant vector a three-cycle gives a nonzero difference of weight
two. Conjugating by $A_m$, which is transitive on unordered pairs,
generates $E$. A family independent under generation by the group
has at most two members, so the normal-module inequality and
$\mu(A_m)\leq m-2$ give $d\leq m$.

The exceptional degree-six $S_5=PGL_2(5)$ is three-transitive.
Every nonconstant vector of weight one, two or three (or its complement)
has conjugate differences generating $E$; for weight three use
three-transitivity, and for the other cases two-transitivity suffices.
Thus it has the same subspace list. Since $m=6$ is even, every invariant
subspace has module-independent length at most one: an odd vector
generates $V$, a nonconstant even vector generates $E$, and a constant
vector generates its line. The normal-module inequality gives
$d\leq\mu(S_5)+1\leq5<6$.

It remains to consider $P=S_m$. The same module list gives
$\mu_G'(W)\leq2$, so $d\leq m+1$. When $m$ is even the preceding
one-generator argument improves this to $d\leq m$.
Suppose $m$ is odd and $d=m+1$. Equality in the normal-quotient proof
forces $|L|=m-1$. Its block images form a Coxeter path: this is
UB.Extremal for $m\geq7$, and the complete independent-string
certificate for $m=5,6$. They occupy an interval in the full string.
Each of the two outside generators centralizes the block subgroup
omitting at most one endpoint of that path. Its centralizer in $S_m$
is trivial (also for the outer $S_6$ image), so both outside generators
actually belong to $V$.

A vector commuting with every adjacent block transposition except
possibly the first or last is constant off that exceptional endpoint.
The possible nonzero vectors are therefore $\lambda,e_i,\lambda+e_i$,
where $i$ is an endpoint of the block path. The vector $e_i$ generates
all of $V$ under $S_m$; two nonconstant even vectors each generate $E$.
Independence of both outside generators consequently forces them to be
$\lambda$ and $\lambda+e_i$. In particular $W=V$.
The subgroup $\gen L$ has zero intersection with $V$: any nonzero
$S_m$-submodule there contains either $\lambda$ or $E$, making one
of these two outside generators redundant. Thus $\gen L$ is a
complement, the graph of a $1$-cocycle $c:S_m\to V$. The automorphism
\[
 (v,\sigma)\longmapsto(v+c(\sigma),\sigma)
\]
of $V\rtimes S_m$ fixes the base module and sends this complement
to the standard one. As an unordered marked set, up to a group isomorphism carrying the
set of distinguished generators onto the other set, the marking is
exactly
\[
 \lambda,\quad\lambda+e_1,\quad s_1,\ldots,s_{m-1},
 \qquad s_i=(i,i+1).
\]
Its displayed ordered Schlaefli type is $[2,4,3,\ldots,3]$.
The noncentral path may be ordered in either direction, and the isolated
central generator may be placed at either end. These are the four
compatible string orders, forming two classes under global duality,
with representatives $[2,4,3,\ldots,3]$ and $[2,3,\ldots,3,4]$.
We do not identify these two ordered polytopes: their labelled types
differ and neither is the global reverse of the other. They share the
same unordered marked set. In particular the source's page-two type
belongs to the reversed noncentral path, not to its literal displayed
tuple; its page-seven type matches that tuple. The source's uniqueness
claim for the unordered pair $(G,S)$ is retained with this distinction.

For completeness the marking has IP for every odd $m\geq3$.
The central factor $\gen\lambda$ is disjoint from $E\rtimes S_m$.
In the remaining family $v=\lambda+e_1,s_1,\ldots,s_{m-1}$, a subset
of adjacent transpositions generates the Young subgroup $P_T$ on
its interval partition. If $v$ is absent the base part is zero.
If $v$ is present and the initial interval containing $1$ is $A$,
the base part is
\[
 W_A=\operatorname{span}\{\lambda+e_i:i\in A\}.
\]
For proper $A$ the indicated vectors are independent, since their
common outside coordinate determines the sum of their coefficients.
For $A=[m]$ their only relation is the sum of all $m$ vectors, giving
$W_{[m]}=E$. All initial intervals are nested and therefore
$W_A\cap W_B=W_{A\cap B}$. The corresponding Young subgroups satisfy
$P_T\cap P_U=P_{T\cap U}$, because their interval partitions refine
by the intersection of their edge sets. The actual subgroups are
$W_A\rtimes P_T$, so their intersections satisfy IP. Adjoining the
independent central factor preserves it. The full group is transitive
on the $2m$ signed points and is $C_2\wr S_m$.

\section{Two general equality lemmas}
\begin{lemma}\label{intransitive}
Let $H\leq S_k$ be intransitive and generated irredundantly by $k-2$
elements. Then $H=S_A\times S_B$ on exactly two nonempty orbits
$A,B$ partitioning the $k$ letters.
\end{lemma}
\begin{proof}
Successive orbit projections and the normal-quotient normalization
bound give independent length at most $\sum_O(|O|-1)=k-q$, where
$q$ is the number of orbits. Hence $q=2$. Write $a=|A|$, $b=|B|$.
Choose a minimal subfamily $T$ generating the first projection. It
has at most $a-1$ members. Normalize all other generators into the
second-coordinate kernel by words in $T$. They remain independent,
so there are at most $b-1$ of them. Equality forces both bounds to
be attained. Whiston's equality clause makes the first image $S_A$
and the subgroup generated by the normalized extra elements $S_B$.
Thus the whole second factor is in $H$, and the full first image
then gives $H=S_A\times S_B$. An orbit of size one has the trivial
symmetric factor and contributes zero throughout.
\end{proof}

We use the following extremal consequence of UB.Extremal: an independent
string marking of $S_b$ with $b-1$ involutions, for $b\geq7$, is a Coxeter
path. Indeed a sign-modified tree marking cannot consist entirely of
involutions: an edge adjacent to the distinguished edge gives an
element of order three. For the ordinary tree, the string condition
makes its line graph a path, so the tree itself is a path. For $b=5,6$
the complete small certificate verifies the same abstract Coxeter-path
conclusion, including the outer-automorphism image in $S_6$. It also
verifies the trivial centralizer of the subgroup omitting an endpoint.

\begin{lemma}\label{smallfactor}
For $a=2$ or $3$, and $b\geq6$, an involutory string C-marking of
$S_a\times S_b$ of rank $a+b-2$ has two nonempty diagram components.
One generates $S_a$, and the other is the graph of a homomorphism
$S_b\to S_a$; this homomorphism is trivial when $a=3$.
\end{lemma}
\begin{proof}
Choose a minimal subfamily $L$ generating the $S_b$ quotient.
Normal-quotient normalization bounds the remaining family by $a-1$,
so $|L|=b-1$. The images of $L$ are an extremal Coxeter path.
Its noncommuting adjacent vertices must be adjacent in the full
string, so $L$ occupies an interval. Every generator outside that
interval centralizes all but possibly one endpoint of this path.
The endpoint-omission subgroup has trivial centralizer in $S_b$;
therefore every outside generator has trivial $S_b$ image and lies
in the small factor. There are exactly $a-1$ of them.

Put $N=\gen{L}\cap S_a$. Their irredundance relative to $\gen{L}$
implies irredundance relative to $N$ under conjugation by $\gen{L}$.
If $a=2$, the one extra generator forces $N=1$ and is central;
its diagram vertex is isolated, and $\gen L$ is a graph over $S_b$.
If $a=3$, any nontrivial subgroup
of $S_3$, together with a family generating $S_3$, needs at most one
member of that family. Thus two independent extras force $N=1$.
The complement $\gen{L}$ is the graph of a homomorphism
$S_b\to S_3$, whose image is trivial or order two, since $A_b$ is
simple and $b\geq5$. If the image has order two, one of the two
extra involutions lies outside it. That involution and its conjugate
by the image generate $S_3$, making the other extra redundant.
Hence the image is trivial. Now the two small-factor involutions
commute with all of $L$, so the diagram is disconnected.
\end{proof}

\section{The uniform two-large-block obstruction}
Let a transitive string C-group $G$ preserve two blocks of size $k$.
Choose a distinguished generator $r$ interchanging the blocks,
nearest to an endpoint of the string among all such generators.
Identify the blocks through $r$, so $r$ is the pure block swap in
$S_k\wr C_2$. Let $C$ be the other distinguished generators commuting
with $r$, and $R$ the remainder. Only its string neighbors can be in
$R$, so $|R|\leq2$.

The centralizer of $r$ in this wreath product is
$\operatorname{diag}(S_k)\times\gen r$. Multiplying each $c\in C$
by $r$ when necessary normalizes it into the diagonal copy of $S_k$.
The projection is injective on $\gen C$, because the intersection
property gives $r\notin\gen C$. Its image $H$ therefore has the
same involutory string C-marking. Moreover
\[
 \gen C=\{(h,h)r^{\chi(h)}:h\in H\}
\]
for a character $\chi:H\to C_2$. Whiston gives $|C|\leq k-1$,
and consequently $d\leq k+2$.

\begin{theorem}\label{twoblocks}
For $k\geq9$, suppose every transitive string C-group of faithful degree
$k$, other than $S_k,A_k$, has rank at most $k/2+1$. Then a transitive
two-block string C-group on $2k$ letters has $d\leq k$.
\end{theorem}
\begin{proof}
If $d=k+2$, then $|C|=k-1$ and $H=S_k$. Its diagram is a Coxeter
path, so $C$ lies on one side of $r$ and its neighbors: $r$ must be
one place from a string endpoint. The outside neighbor centralizes
all of $C$. But $\ker\chi$ contains $A_k$, whose centralizer in
$S_k$ is trivial. An element $(p,q)r^\epsilon$ centralizing $C$
therefore has $p=q=1$, and is in $\gen r$, a contradiction.

Suppose $d=k+1$. If $r$ is at an endpoint, $|C|=k-1$ and the one
remaining generator centralizes a Coxeter endpoint-omission subgroup
of $C$. Its character kernel contains an $A_{k-1}$ fixing one letter;
that subgroup again has trivial centralizer in $S_k$. The remaining
generator is in $\gen r$, a contradiction. If $r$ is interior but
$|R|=1$, then again $|C|=k-1$ and its diagram is connected. Therefore
$r$ is one place from an endpoint and the sole noncommuting neighbor
is on that outer side. It commutes with all of $C$, so the preceding
$A_k$ centralizer argument excludes it. Thus $|R|=2$ and $|C|=k-2$.

First let $r$ have distance one from an endpoint, so $C$ occupies
one interval (its noncommuting diagram need not be connected).
Its outside neighbor centralizes all of $C$ and does not commute with
$r$. If $H$ is transitive, the hypothesis at degree $k$ and
$k-2>k/2+1$ leave only $S_k,A_k$. This use is legitimate because
the projection onto $H$ preserves the full marking and IP. The same
character-kernel centralizer argument excludes both.
Otherwise Lemma~\ref{intransitive} gives $H=S_A\times S_B$.
If both orbits have size at least four, $\ker\chi$ contains
$A_A\times A_B$ and has trivial centralizer. A singleton orbit is
also excluded by this argument on its complement. The only remaining
possibilities have $a=|A|\in\{2,3\}$ and $b=|B|\geq6$.
If $\chi$ is nontrivial on the large factor, a centralizing
$(p,q)r^\epsilon$ first has $p,q$ supported on $A$ (by $A_B$),
then has $p=q$ by an odd permutation supported on $B$. It commutes
with $r$, a contradiction. If $\chi$ is trivial there,
Lemma~\ref{smallfactor} gives two components. The component at the far
endpoint has extremal Coxeter generators, all odd in its symmetric
factor. A nontrivial restriction of $\chi$ to that component would
make its endpoint a block swap closer to an endpoint than $r$.
Hence $\chi$ is trivial on that far component. Both neighbors of $r$
centralize it: the inner neighbor commutes with all of $C$ except
possibly its first vertex, in the other component. The far component
has $A$ or $B$ as its distinguished orbit; any other nontrivial orbit
has a different size. Its centralizer in $S_k$ therefore preserves
$A,B$. Thus the coordinate permutations of both neighbors preserve
$A,B$. Together with $C$ and $r$ they preserve $A\times\{0,1\}$,
contradicting transitivity.

Finally let $r$ have distance at least two from both endpoints.
Then $C$ occupies two nonempty intervals separated by the missing
neighbors and $r$. Their generated images commute and form a direct
product decomposition of $H$, by the intersection property; we do not
yet assume that either interval has connected noncommuting diagram.
The same smaller-degree hypothesis excludes transitive $H$ other than
$S_k,A_k$, and these cannot be a nontrivial direct product. Thus again
$H=S_A\times S_B$ by Lemma~\ref{intransitive}.
Neither orbit is a singleton. For orbit sizes at least three, the
symmetric factors are centerless and directly indecomposable, so
the two interval factors are precisely those symmetric factors. Their
noncommuting diagrams are therefore connected. If one orbit has
size two, one component is this central $C_2$, and the other is a
graph of a character $S_B\to C_2$. This graph is directly
indecomposable too. Thus in both cases the intervals are exactly the
two noncommuting components, and each attains the maximum independent
length in its symmetric factor.

If $\chi$ is nontrivial on such a component of size at least five,
its extremal Coxeter generators are all odd. Hence the endpoint of
that component at the end of the whole string interchanges the two
large blocks, contradicting the choice of $r$. The same holds for
factors of size two or three. For $S_4$, at least one of the first
two and one of the last two generators is odd: otherwise two even
involutions lie in its abelian normal four-group, and the string
relations make an endpoint generator central in $S_4$.
This again gives a closer block-swapping generator. Therefore
$\chi$ is trivial on both components, and all of $C$ fixes both
large blocks individually.

Each member of $R$ centralizes the opposite component of $C$.
Its coordinate permutations consequently preserve the two sets
$A,B$. This remains true for the possible character graph over the
large symmetric factor when $|A|=2$: its two nontrivial orbits have
different sizes. All of $C$, $R$ and $r$ therefore preserve
$A\times\{0,1\}$ and $B\times\{0,1\}$, contradicting transitivity.
\end{proof}

\section{The finite two-block boundary}
The same argument excludes $d\geq k+2$ for $5\leq k\leq8$.
For $k=5,6$ the complete independent-string certificate proves both
that the extremal markings are Coxeter paths and that the derived
subgroup of an endpoint-omission subgroup has trivial centralizer.
In the outer $S_6$ marking that subgroup is a transitive degree-six
$A_5$, not a point-fixing alternating group. These actual centralizer
checks replace the natural point-fixing description in the endpoint
argument. The cases with $d=k+1$ and $|R|\leq1$ are excluded in the
same way.

For $d=k+1$ and $|R|=2$, the faithful normalized $H\leq S_k$ has
rank $k-2$. If it is intransitive, Lemma~\ref{intransitive} gives
$H=S_a\times S_b$. A singleton orbit is excluded by the alternating
centralizer argument, or by direct indecomposability when the two
intervals are nonempty. The remaining unordered orbit sizes are
$(2,3),(2,4),(3,3),(2,5),(3,4),(2,6),(3,5),(4,4)$.
For a transitive $H$, the accepted primitive bounds and the complete
small transitive-imprimitive census leave only $S_k,A_k$, the
primitive degree-six $S_5$, and the degree-six actions numbered
9 and 11. The first two are excluded as above. In the degree-six
$S_5$ case, its derived subgroup is primitive $A_5$ with trivial
centralizer, so the same argument applies. The degree-eight census
has no rank-six marking; degree five and seven have no nontrivial
imprimitive actions.

For the eight product actions and the two degree-six actions, the
separate, completely proved finite-neighbor lemma enumerates every
$C$-marking, character, compatible split and pair of missing neighbors.
Its 59 normalized markings yield 154 transitive completions, all
redundant. This exact required finite premise therefore excludes all
remaining cases. Consequently the two-large-block bound $d\leq k$
holds for every $k\geq5$. The block sizes three and four are included
in the separately complete degree-six and degree-eight census; they
contain genuine exceptional markings.

\section{A sharp primitive bound with a bounded residual census}
We now assume that $G$ is primitive of degree $n$, distinct from $A_n,S_n$.
If its noncommuting diagram is disconnected, the earlier primitive
commuting-factor argument gives $G=T\times T$ on $|T|=n\geq60$ points
and $d\leq2\log_2n<n/2$. Otherwise every adjacent pair fails to commute.
We explicitly retain the following independent theorem of Conder:
a finite nondegenerate regular polytope of rank $r\geq9$ has at least
$2^{2r-1}$ flags, equivalently its string C-group has at least that order.
We use only this stated rank range and only connected string diagrams.
This is an external theorem, not a consequence of our finite diagnostics.

Apply Mar\'oti's original primitive-order theorem, including its lower
containment $A_m^\ell\leq G\leq S_m\wr S_\ell$ in the product-action
alternative, with $m\geq5$ and, after replacing subsets by complements,
$1\leq k\leq m/2$. The degree is $n=\binom m k^\ell$.
If $\ell=1$, the natural $k=1$ action would give the excluded $A_n,S_n$.
For $k\geq2$, Whiston in the faithful $m$-point action gives
$d\leq m-1<n/2$. If $\ell\geq2$, the same abstract group has a faithful
transitive imprimitive action on $m\ell$ letters: the lower subgroup
$A_m^\ell$ is transitive in each block and primitivity of the product
action makes the coordinate action transitive. The induction hypothesis at the strictly smaller degree $m\ell<n$
provides the imprimitive bound and gives $d\leq m\ell/2+1<n/2$, since
$m\ell+2<m^\ell\leq\binom m k^\ell$ for $m\geq5,\ell\geq2$.

The Mathieu alternative is treated by the bounded checks below.
In the remaining alternative put
\[
 M(n)=n\prod_{i=0}^{\lfloor\log_2n\rfloor-1}(n-2^i).
\]
Exact integer arithmetic verifies $M(n)<2^{n-1}$ for $26\leq n\leq63$.
For $n\geq64$, with $q=\lfloor\log_2n\rfloor\geq6$, we have
$M(n)<n^{q+1}<2^{(q+1)^2}\leq2^{n-1}$; the inequality
$(q+1)^2\leq2^q-1$ holds at $q=6$ and persists by induction.
If $d\geq n/2$, then $d\geq13$, so Conder gives
$|G|\geq2^{2d-1}\geq2^{n-1}>M(n)$, a contradiction.
We use the exact product bound here, not the source's coarser logarithmic
estimate at its asserted small threshold.

\subsection{Complete finite screening}
For $2\leq n\leq25$, the required program visits every primitive
catalogue entry, omitting exactly $A_n,S_n$. It uses the following
necessary bounds, with all finite groups and actions computed exactly.
If $\Omega(t)$ counts prime factors with multiplicity, then
$d\leq\Omega(|G|)$, by the strictly increasing subgroup chain from IP.
All distinguished generators are involutions, so a proper subgroup
generated by all involutions excludes $G$ outright. For every possible
endpoint $x$, the subgroup on labels $2,\ldots,d-1$ commutes with $x$
and intersects $\gen x$ trivially. Consequently
\[
 d\leq1+\Omega(|C_G(x)|).
\]
Likewise the even labels generate an elementary abelian subgroup of
rank $\lceil d/2\rceil$. If every elementary abelian subgroup of a
Sylow two-subgroup has rank at most $e$, then $d\leq2e$.
The finite program computes this bound by the complete subgroup list
of that Sylow group; no assertion about its mere minimum number of
generators is substituted for its elementary abelian rank.

For primitive affine degree sixteen, the program verifies the regular
normal subgroup $2^4$. The normal-module lemma costs one generator,
and the faithful quotient embeds in $GL_4(2)\cong A_8$, whose independent
families have length at most six by Whiston. Thus $d\leq7<8$.
The exceptional isomorphism is explicitly checked and is also the
retained identification already used by common Topic B.
For the degree-nine affine group $AGL_2(3)$, the program verifies
$\Phi(GL_2(3))\cong C_2$ and the Frattini quotient $S_4$; its maximum
irredundant generating length is at most three. The irreducible
translation module costs one, giving $d\leq4<9/2$.
The degree-fifteen copy of $A_8$ has $d\leq6$ directly.

For both actions of $M_{11}$, the Sylow elementary abelian rank is two,
so $d\leq4$. The degree-twenty-three group $M_{23}$ is already excluded
by the endpoint-centralizer bound. For $M_{24}$ the two involution
centralizers have orders $7680,21504$. If $d\geq12$, the connected tail
on labels $2,\ldots,d-1$ has rank at least ten. Conder applied to that
tail, disjoint from $\gen x$, requires
$|C_G(x)|\geq2\cdot2^{2(d-2)-1}=2^{2d-4}\geq2^{20}$,
contradicting both actual centralizer orders.

The surviving actions are exactly the following GAP entries:
\[
 (6,1),(6,2),(7,5),(8,3),(8,5),(10,4),(10,5),(10,7),(12,2).
\]
The degree-eight copy $(8,4)$ is explicitly identified with $(7,5)$.
These are GAP's current indices, not Sims' indices printed in the source.
The maximum ranks searched are respectively $3,4,4,5,4,6,6,6,6$.
For $A_5,S_5$ they follow from Whiston; for $PSL_3(2)$ and $AGL_3(2)$
from the accepted common independent-length bounds; for the others
from the checked Sylow elementary abelian ranks two or three.

We enumerate connected involutory strings up to simultaneous normalizer
conjugation. At each prefix we discard dependence, failed string
commutations, or a failed endpoint-parabolic intersection. These are
necessary conditions for every subfamily of a C-marking. If $H$ is the
prefix subgroup and at least $t$ future generators remain after the next
one, their group centralizes $H$ and is disjoint from it; hence
$[C_G(H):Z(H)]\geq2^t$. This necessary inequality supplies a further
pruning rule. At a full generating tuple every pair of subset subgroups
is tested for IP. Stabilizer orbits at each extension preserve the full
prefix and cannot discard an ordered conjugacy class. No sample
markings, probabilistic generation test or source rank table is used.

The complete outputs, up to duality, are
\[
\begin{array}{c|c|l}
 n&G&\text{types at rank at least }n/2\\\hline
6&A_5&[3,5],\ [5,5]\\
6&S_5&[3,3,3],\ [4,5],\ [4,6],\ [5,6],\ [6,6]\\
10&S_6&[3,3,3,3].
\end{array}
\]
The last group is GAP $(10,5)$ and is identified exactly with $S_6$.
The full permutation tuples are retained, including both dual orientations.
Only the rank-four degree-six simplex exceeds half the permutation degree.
Thus this recovers the complete primitive branch of Theorem~1.2 and the
exception list of Proposition~3.2. The source's finer Table~3 maximum-rank
numbers are not imported as premises of this proof; any not independently
computed are separately attributed ancillary reports in the source map.

\section{Closing the simultaneous induction}
\begin{theorem}\label{sharpclosure}
For every faithful transitive string C-group of degree $n$, other than
$A_n,S_n$, the rank is at most $n/2+1$. In the imprimitive case with
$n\geq10$, equality requires two-point blocks with an odd number of
blocks, and has exactly the markings classified above. In the primitive
case the rank is strictly below $n/2$ except for the explicitly enumerated
degree-six and degree-ten markings.
\end{theorem}
\begin{proof}
We prove these assertions simultaneously by strong induction on $n$.
The small imprimitive actions of degrees six, eight and nine are covered
by the complete required census. It searches through rank $n-2$, the
Whiston bound for a proper subgroup other than $S_n$, without using any
imprimitive rank theorem. A prime degree has no nontrivial block system.
For degree four the possible imprimitive groups are subgroups of $D_8$;
its Frattini quotient is $C_2^2$, so any independent generating family
has at most two members. The primitive checks through degree twenty-five
are independent of the desired induction assertion. These observations
supply every small endpoint needed below.

At degree $n$, the maximal-block argument handles $k,m>2$. For blocks
of size two it uses only the previously accepted independent-length
bounds on primitive quotient images, the explicit binary-module lemma,
and the complete extremal $S_5,S_6$ checks. It does not apply a C-group
rank theorem to a quotient that might have lost IP. For two blocks, the
finite ten-action lemma treats $5\leq k\leq8$. When $k\geq9$, the
normalized $H$ is itself a string C-group of degree $k<n$, so the
induction hypothesis supplies precisely the assumption in
Theorem~\ref{twoblocks}. Thus the whole imprimitive conclusion at degree
$n$ follows with its complete equality classification.

For a primitive group, the product-action alternative of Mar\'oti has
a faithful imprimitive action of degree $m\ell<n$. The induction
hypothesis, not a same-degree result, applies there. The other primitive
alternatives use Whiston, the retained Conder theorem, or the independent
bounded catalogue checks, as proved above. This proves the primitive
conclusion at degree $n$. Every invocation of the asserted sharp rank
bound has therefore decreased the faithful permutation degree, closing
the induction without using the older general imprimitive
independent-length theorem.
\end{proof}

\appendix
\section{Complete retained finite-neighbor proof}
The following full body is copied exactly from the root-authored bounded note,
independently accepted by the separate finite-boundary reviewer. Its original
wrapper is retained in the archive, and its entire mathematical content is
part of this common manuscript. The complete programs appear in our appendix.
We establish the finite obstruction used in the two-block argument.
The classification of the possible groups in this statement is a
separate obligation: no conclusion about other groups is implicit in
the calculation below. In particular, this component is not a proof
of the full transitive rank theorem.

Let $W=S_k\wr C_2$ act on two blocks of size $k$, and let $r$ be the
permutation exchanging their corresponding letters. For a subgroup
$H\leq S_k$ and a character $\chi:H\to C_2$, put
\[
 D_\chi(H)=\{(h,h)r^{\chi(h)}:h\in H\}.
\]
The two symmetric factors in a row $S_a\times S_b$ of the table below
act on disjoint sets of sizes $a,b$ in one block. The two transitive
degree-six actions are the literal permutation groups numbered 9
and 11 in the GAP transitive-group catalogue; their orders are 36
and 48, respectively. The action, rather than the abstract group
alone, is part of the specification.

\begin{lemma}\label{boundary}
Let $H$ be one of the ten actions in Table~1, of degree $k$.
Take an involutory string C-marking $(h_1,\ldots,h_{k-2})$
of $H$, and let $C$ be its image in $D_\chi(H)$. Split this ordered
list as $C_L,C_R$, requiring the two lists to commute elementwise.
There is no tuple
\[
                  (C_L,x,r,y,C_R)
\]
of independent involutions that generates a transitive subgroup of
$W$, satisfies the string commutation relations, has both $x$ and
$y$ noncommuting with $r$, and has $r$ at least as close to an
endpoint as every other generator exchanging the blocks.
\end{lemma}

\begin{proof}
We give the complete finite enumeration and its exhaustion argument.
The accompanying program uses exact finite permutation groups.

First enumerate the involutory string C-markings of $H$. At a prefix
$(h_1,\ldots,h_j)$, an eligible next generator is an involution of $H$
outside the prefix subgroup and commuting with $h_1,\ldots,h_{j-1}$.
We discard a prefix if one of its generators belongs to the subgroup
generated by the others. Every subfamily of an independent family
is independent, so this discards no desired tuple. Also, if the
current subgroup has index less than $2^{k-2-j}$ in $H$, the prefix
cannot acquire the required number of independent generators: each
strict subgroup enlargement has index at least two.

Simultaneous conjugation by $N_{S_k}(H)$ preserves every condition.
At each step choose one orbit representative under the pointwise
stabilizer of the prefix under conjugation, then replace that
stabilizer by its centralizer of the chosen generator. Induction on
the prefix length shows that every ordered marking is represented.
At length $k-2$, require generation of $H$ and test the intersection
property for every pair of subsets of the marking. The test forms
the actual generated subgroups and their intersections, without a
quotient or an assumed presentation.

For each surviving marking enumerate all $2^{k-2}$ assignments of
its generators to $C_2$. Retain exactly those assignments for which
\texttt{GroupHomomorphismByImages} constructs a homomorphism.
Since the marking generates $H$, these are exactly its characters.
Construct the corresponding diagonal lifts explicitly as
permutations of $2k$ letters. Try every cut between the lifts,
including the two empty-end cuts, and require elementwise
commutation across it. Reject a cut and character only if an
already fixed block-exchanging generator is closer to an endpoint
than $r$.

The string relations require $x$ to centralize all of $C_R$ and
all but the last member of $C_L$. Similarly, $y$ must centralize
all of $C_L$ and all but the first member of $C_R$. Compute these
two centralizers inside $W$, enumerate their involutions, and impose
noncommutation with $r$ and the same endpoint restriction. We may
choose one representative for $x$ under the simultaneous centralizer
of $C$ and $r$ in $W$: conjugation by this group preserves the complete
candidate set for $y$. For each remaining pair require $xy=yx$ and
transitivity. Finally test independence by deleting each generator
in turn and comparing the resulting subgroup with the full group.

The exhaustive counts are in Table~1. No transitive completion occurs
in the first eight cases. In the last two cases, all transitive
completions fail independence. Thus no test of the full intersection
property is needed to reject a completion, proving the lemma.
\end{proof}

\begin{table}[ht]
\centering
\begin{tabular}{lrrrr}
\toprule
Action of $H$ & $k$ & Markings & Transitive completions & Independent\\
\midrule
$S_2\times S_3$ &5&4&0&0\\
$S_2\times S_4$ &6&12&0&0\\
$S_3\times S_3$ &6&1&0&0\\
$S_2\times S_5$ &7&4&0&0\\
$S_3\times S_4$ &7&6&0&0\\
$S_2\times S_6$ &8&8&0&0\\
$S_3\times S_5$ &8&2&0&0\\
$S_4\times S_4$ &8&9&0&0\\
$\mathrm{TransitiveGroup}(6,9)$ &6&1&2&0\\
$\mathrm{TransitiveGroup}(6,11)$ &6&12&152&0\\
\bottomrule
\end{tabular}
\caption{Complete bounded search. Markings are represented up to the
normalizer conjugations described in the proof. Completion counts
use the stated pruning of $x$; they are not counts of isomorphism classes.}
\end{table}

The implementation is in \texttt{search.g}; \texttt{run.g} executes
all ten cases, and \texttt{check.py} requires their exact summaries
and successful completion. To check that the procedure can retain
valid completions, it is also run with $k=3$ and the empty marking
of the trivial group. It recovers six independent transitive
rank-three string C-markings after the same pruning. For example,
\[
       ((5\ 6),\ (1\ 4)(2\ 5)(3\ 6),\ (1\ 2))
\]
passes the full intersection test. A redundant triple of distinct
transpositions in $S_3$ supplies negative controls for independence
and the intersection property.

The proof depends on GAP's exact finite permutation-group algorithms
and its degree-six transitive catalogue. It does not supply a Lean
formalization. The separate reduction of an arbitrary two-block
group to these ten actions, the endpoint cases, and all other
branches of the transitive rank theorem remain outside this lemma.
Normalizing the commuting generators by powers of the block swap
already appears in Cameron, Fernandes, Leemans and Mixer,
\emph{String C-groups as transitive subgroups of $S_n$},
arXiv:1410.5863v1 (2014), proof of Proposition~2.1, Case~3.
We retain that step. The present contribution is the explicit
character description and exhaustive finite analysis with two
noncommuting neighbors, within the common string-separator argument.
It neither claims a new open-conjecture solution nor cites the
target rank theorem as a premise.

\begin{thebibliography}{9}
\bibitem{CFLM14} P. J. Cameron, M. E. Fernandes, D. Leemans and M. Mixer,
\emph{String C-groups as transitive subgroups of $S_n$},
arXiv:1410.5863v1, 21 October 2014, twelve-page edition.
\bibitem{CC02} P. J. Cameron and P. Cara, \emph{Independent generating
sets and geometries for symmetric groups}, J. Algebra 258 (2002), 641--650.
\bibitem{W00} J. Whiston, \emph{Maximal independent generating sets of
the symmetric group}, J. Algebra 232 (2000), 255--268.
\bibitem{M02} A. Mar\'oti, \emph{On the orders of primitive groups},
J. Algebra 258 (2002), 631--640, Theorem 1.1 in its original form.
\bibitem{C13} M. Conder, \emph{The smallest regular polytopes of given
rank}, Adv. Math. 236 (2013), 92--110. The nondegenerate rank-at-least-nine
lower bound is retained as stated in the text; its full proof is not supplied here.
\end{thebibliography}
\end{document}
